% ======================== 第十一节 声明（签字页） ============================
% 版式与结构对照 2026 年真实申报稿：各方声明用「一、二、…」编号标题（进目录），
% 签名区块＝左列标签＋右区签名线阵列（\sigblock + \sigzrow，单人签名位靠右），
% 单位盖章与年月日右下落款（\sigseal，电子章图片钩子见 zhaogushu.sty 注释）。
% 各方声明措辞照准则第 57 号第八十五至九十条；治理结构按新《公司法》
% （不设监事会），签字栏分「全体董事／全体审计委员会成员／除董事、审计委员会
% 成员以外的高级管理人员」。
\gssection{声明}

\subsection{发行人及全体董事、审计委员会成员、高级管理人员声明}

本公司及全体董事、审计委员会成员、高级管理人员承诺本招股说明书的内容真实、
准确、完整，不存在虚假记载、误导性陈述或重大遗漏，按照诚信原则履行承诺，并
承担相应的法律责任。

\sigblock{全体董事：}{%
\sigzrow{林××}{林×杰}{林×桐}
\sigzrow{林×阳}{林×雪}{王××}
\sigzrow{赵××}{钱××}{孙××}}

\sigblock{全体审计委员会成员：}{%
\sigzrow{赵××}{钱××}{孙××}}

\sigline{除董事、审计委员会成员以外的高级管理人员：}{卫××}{}

\sigseal{发行人（盖章）}{\gsCompany}

\clearpage
\subsection{发行人控股股东、实际控制人声明}

本公司及本人承诺本招股说明书的内容真实、准确、完整，不存在虚假记载、误导性
陈述或重大遗漏，按照诚信原则履行承诺，并承担相应的法律责任。

\sigline{控股股东法定代表人：}{林××}{}

\sigline{实际控制人：}{林××}{}

\sigseal{控股股东（盖章）}{示例控股有限公司}

\clearpage
\subsection{保荐人（主承销商）声明}

本公司已对招股说明书进行核查，确认招股说明书的内容真实、准确、完整，不存在
虚假记载、误导性陈述或重大遗漏，并承担相应的法律责任。

\sigline{保荐代表人：}{徐××}{何××}

\sigline{项目协办人：}{郭××}{}

\sigline{法定代表人：}{吕××}{}

\sigseal{保荐人（主承销商）（盖章）}{\gsSponsor}

\clearpage
\subsection{发行人律师声明}

本所及经办律师已阅读招股说明书，确认招股说明书与本所出具的法律意见书无矛盾
之处。本所及经办律师对发行人在招股说明书中引用的法律意见书的内容无异议，确
认招股说明书不致因上述内容而出现虚假记载、误导性陈述或重大遗漏，并承担相应
的法律责任。

\sigline{经办律师：}{谢××}{董××}

\sigline{律师事务所负责人：}{袁××}{}

\sigseal{律师事务所（盖章）}{示例市示例律师事务所}

\clearpage
\subsection{会计师事务所声明}

本所及签字注册会计师已阅读招股说明书，确认招股说明书与本所出具的《审计报告》、
《审阅报告》、《内部控制鉴证报告》及经本所鉴证的非经常性损益明细表无矛盾之
处。本所及签字注册会计师对发行人在招股说明书中引用的上述报告及文件的内容无
异议，确认招股说明书不致因上述内容而出现虚假记载、误导性陈述或重大遗漏，并
承担相应的法律责任。

\sigline{签字注册会计师：}{高××}{秦××}

\sigline{会计师事务所负责人：}{罗××}{}

\sigseal{会计师事务所（盖章）}{示例会计师事务所（特殊普通合伙）}

% --- 资产评估机构声明（如本次发行涉及评估事项，取消注释并补全；措辞照准则
%     第 57 号第九十条）------------------------------------------------------
% \clearpage
% \subsection{资产评估机构声明}
%
% 本机构及签字资产评估师已阅读招股说明书，确认招股说明书与本机构出具的资产
% 评估报告无矛盾之处。本机构及签字资产评估师对发行人在招股说明书中引用的资
% 产评估报告的内容无异议，确认招股说明书不致因上述内容而出现虚假记载、误导
% 性陈述或重大遗漏，并承担相应的法律责任。
%
% \sigline{签字资产评估师：}{××}{××}
%
% \sigline{资产评估机构负责人：}{××}{}
%
% \sigseal{资产评估机构（盖章）}{示例资产评估有限公司}
